%==============================================================================
%  MATH 231 -- Linear Algebra I (Queens College, Fall 2026)
%  MATRIX UNIT, PART 8 -- Solving Linear Systems: Elimination
%
%  Section numbering runs CONTINUOUSLY across the parts of the matrix unit;
%  see Part 2's header for the full map.  This part is Sections 40-45, hence
%  \setcounter{section}{39}.
%
%  STATUS: complete.  Sections 40-45.  Written from the plan in
%  plan/MATH231_remaining_units_plan.tex (2026-10-25).
%
%  AUDIENCE: distributed to students.  No instructor-facing apparatus.
%
%  WHAT THE INSTRUCTOR SPECIFIED, AND WHERE IT IS.
%   - n x n systems, Gaussian elimination, back/forward substitution
%     (Sec. 41-42); Gauss-Jordan elimination (Sec. 43); row-echelon form
%     (Sec. 42.5) and reduced row-echelon form (Sec. 43.1); elementary row
%     operations (Sec. 40); homogeneous systems (Sec. 44); overdetermined and
%     underdetermined systems (Sec. 45).
%
%  DEBTS COLLECTED.
%   - Part 1's closing Looking-ahead (end of Sec. 9) promised Gaussian
%     elimination, and Sec. 4.4 hinted that it makes the matrix triangular.
%     Sec. 42.2 solves Part 1's own 3x3 system (Sec. 1.6), recovering
%     alpha = [1, 2, 3].
%   - D5 (plan): the second, computational proof of rank-nullity by counting
%     pivot and free columns is Sec. 43.4.
%   - D6/D10 (plan; the determinant law D10 of Sec. 38): "rank n => det != 0"
%     is completed in Sec. 42.6 by elimination.
%
%  PSEUDOCODE uses the algo environment copied from computational unit
%  Part 4, so the two units' algorithms look the same.
%
%  OPERATION COUNTS are given in round terms (n^2, (2/3) n^3) with a short
%  justification, without the FLOP bookkeeping of computational unit Part 1.
%
%  THE LINE-FIT DATA of Sec. 45.1 -- (1,1), (2,2), (3,2) -- is reused in
%  Part 9 (Sec. 49.4) and Part 11 (Sec. 58.2).  Keep them in step.
%
%  Compile with pdflatex or tectonic (standard TeX Live packages only).
%==============================================================================
\documentclass[11pt]{article}

\usepackage[margin=1in]{geometry}
\usepackage{amsmath,amssymb,amsthm}
\usepackage[dvipsnames]{xcolor}
\usepackage{enumitem}
\usepackage{booktabs}
\usepackage{array}
\usepackage[hidelinks]{hyperref}
\usepackage{titlesec}
\usepackage{needspace}

%------------------------------------------------------------------ style ----
\titleformat{\section}{\large\bfseries\sffamily}{\thesection.}{0.6em}{}
\titleformat{\subsection}{\normalsize\bfseries\sffamily}{\thesubsection.}{0.5em}{}
\titlespacing*{\section}{0pt}{14pt}{6pt}
\setlength{\emergencystretch}{3em}

\theoremstyle{plain}
\newtheorem{thm}{Theorem}[section]
\newtheorem{lem}[thm]{Lemma}
\newtheorem{cor}[thm]{Corollary}

\theoremstyle{definition}
\newtheorem{defn}{Definition}[section]
\newtheorem{ex}{Example}[section]
\theoremstyle{remark}
\newtheorem*{rmk}{Remark}
\newtheorem*{warn}{A common confusion}

\newcommand{\lookahead}{\par\smallskip\noindent
  \textbf{\sffamily\textcolor{RoyalBlue}{Looking ahead.}}\ \ignorespaces}
\newcommand{\lookback}[1][the vector unit]{\par\smallskip\noindent
  \textbf{\sffamily\textcolor{ForestGreen}{From #1.}}\ \ignorespaces}

%--------------------------------------------------------------- notation ----
\newcommand{\R}{\mathbb{R}}
\newcommand{\vc}[1]{\mathbf{#1}}
\newcommand{\mt}[1]{\mathbf{#1}}
\newcommand{\norm}[1]{\lVert #1\rVert}
\newcommand{\T}{^{\mathsf{T}}}
\newcommand{\colsp}[1]{\operatorname{col}[#1]}
\newcommand{\nullsp}[1]{\operatorname{null}[#1]}
\DeclareMathOperator{\spanop}{span}
\DeclareMathOperator{\diagop}{diag}
\DeclareMathOperator{\rank}{rank}
\DeclareMathOperator{\nullity}{nullity}
% an augmented matrix [A | b]: \begin{aug}{ccc|c} ... \end{aug}
\newenvironment{aug}[1]{\left[\begin{array}{#1}}{\end{array}\right]}

%--------------------------------------------------------------- pseudocode --
\newcommand{\kw}[1]{\textsf{\bfseries #1}}
\newcommand{\ind}{\hspace*{1.3em}}
\newcommand{\indd}{\hspace*{2.6em}}
\newcounter{algoline}
\newenvironment{algo}[2]{%
  \par\medskip\needspace{8\baselineskip}%
  \begin{center}\begin{minipage}{0.94\textwidth}%
  \hrule height 0.9pt \vspace{3pt}%
  \noindent{\sffamily\bfseries #1}\par\vspace{2pt}%
  \noindent{\small #2}\par\vspace{3pt}%
  \hrule height 0.4pt \vspace{5pt}%
  \begin{list}{\textbf{\arabic{algoline}:}}%
    {\usecounter{algoline}\setlength{\leftmargin}{2.3em}%
     \setlength{\labelwidth}{1.8em}\setlength{\labelsep}{0.5em}%
     \setlength{\itemsep}{2pt}\setlength{\topsep}{0pt}%
     \setlength{\parsep}{0pt}}%
}{%
  \end{list}\vspace{5pt}\hrule height 0.9pt%
  \end{minipage}\end{center}\medskip%
}

\title{\vspace{-1.2cm}\textbf{\sffamily MATH 231 -- Linear Algebra I}\\[4pt]
       {\large\sffamily The Matrix Unit, Part 8 \textbullet\ Solving Linear
        Systems: Elimination}\\[2pt]
       {\normalsize\sffamily Kiely Hall 326 \textbullet\
        Instructor: Kevin Bedoya}}
\author{}
\date{}

\begin{document}
\setcounter{section}{39}  % this part opens at Section 40; see the header note
\maketitle
\vspace{-1.6cm}

\noindent\textbf{About this document.} This is Part 8 of the matrix unit,
Sections 40--45. Section numbers run continuously through the unit (Part 2's
opening table lists them all).

\medskip
\noindent\textbf{Assumed.} Linear systems and $\mt{A}\vc{x}=\vc{b}$ (Part 1);
triangular matrices (\S4.4); column space, null space, rank and rank--nullity
(\S10--\S12, \S31--\S32); the determinant and its laws D1--D10 (\S38--\S39).

\medskip
\noindent\textbf{What you should be able to do afterwards.} Perform the three
elementary row operations and write each as a matrix; solve triangular systems
by back and forward substitution; reduce a matrix to row-echelon form by Gaussian
elimination and solve the system; continue to reduced row-echelon form by
Gauss--Jordan elimination; read off the rank, the pivot and free variables, a
basis of the null space, and every solution; decide from the reduced form whether
a system has no solution, one, or infinitely many; compute a determinant by
elimination; and describe what to expect of over- and underdetermined systems.

\bigskip
\noindent Part 1 wrote a linear system as $\mt{V}\vc{x}=\vc{u}$, solved a
$2\times2$ example by substitution, and noted that substitution is not a method
that survives twenty equations. It promised a procedure that works on the matrix
directly, applies to any size, and ends either with the answer or with a definite
report that there is none. Since then, Parts 2--7 have explained what the
possible answers are. This part supplies the procedure.

%==============================================================================
\section{Elementary row operations}
%==============================================================================

%------------------------------------------------------------------------------
\subsection{The augmented matrix}
%------------------------------------------------------------------------------
The system $\mt{A}\vc{x}=\vc{b}$ is completely described by the coefficients and
the right-hand side. Writing $\vc{b}$ as an extra column, separated by a bar, gives
the \textbf{augmented matrix} $[\,\mt{A}\mid\vc{b}\,]$. Each row is one equation:
\[
  \begin{aligned}
    x_1+2x_2\phantom{{}+x_3}&=5\\
    x_2+x_3&=5\\
    x_1\phantom{{}+2x_2}+x_3&=4
  \end{aligned}
  \qquad\longleftrightarrow\qquad
  \begin{aug}{ccc|c}1&2&0&5\\0&1&1&5\\1&0&1&4\end{aug}.
\]
This is the system of Part 1 \S1.6, whose solution was constructed to be
$\vc{x}=[\,1,2,3\,]$. Solving it, by a method, is \S42's first job.

%------------------------------------------------------------------------------
\subsection{Three operations}
%------------------------------------------------------------------------------
\begin{defn}[Elementary row operations]
The three \textbf{elementary row operations} on a matrix are:
\begin{enumerate}[label=\textbf{(E\arabic*)},leftmargin=3em,itemsep=2pt]
  \item \textbf{swap} two rows;
  \item \textbf{scale} a row by a nonzero number;
  \item \textbf{replace} a row by itself plus a multiple of another row.
\end{enumerate}
\end{defn}

\noindent On the augmented matrix these are things one does to equations without
thinking: list them in a different order, multiply one through by a nonzero
number, add a multiple of one equation to another.

%------------------------------------------------------------------------------
\subsection{They do not change the solutions}
%------------------------------------------------------------------------------
\begin{thm}
\label{thm:rowops}
If $[\,\mt{A}'\mid\vc{b}'\,]$ is obtained from $[\,\mt{A}\mid\vc{b}\,]$ by an
elementary row operation, the systems $\mt{A}\vc{x}=\vc{b}$ and
$\mt{A}'\vc{x}=\vc{b}'$ have exactly the same solutions.
\end{thm}

\begin{proof}
Every equation of the new system is a combination of equations of the old one,
so any solution of the old system satisfies the new one. And every operation can
be undone by an operation of the same kind: swap the same rows back; scale by the
reciprocal; subtract the multiple that was added. So the old system is also
obtained from the new one, and any solution of the new system satisfies the old
one.
\end{proof}

%------------------------------------------------------------------------------
\subsection{Each operation is a matrix}
%------------------------------------------------------------------------------
Doing a row operation to $\mt{A}$ is the same as multiplying $\mt{A}$ on the
left by the matrix $\mt{E}$ obtained by doing that operation to $\mt{I}$.
Row $i$ of $\mt{E}\mt{A}$ is (row $i$ of $\mt{E}$) times $\mt{A}$, a combination
of the rows of $\mt{A}$ with coefficients from row $i$ of $\mt{E}$. In $\R^{3\times3}$:
\[
  \underbrace{\begin{bmatrix}0&1&0\\1&0&0\\0&0&1\end{bmatrix}}_{\text{swap rows 1, 2}},
  \qquad
  \underbrace{\begin{bmatrix}1&0&0\\0&c&0\\0&0&1\end{bmatrix}}_{\text{scale row 2 by }c},
  \qquad
  \underbrace{\begin{bmatrix}1&0&0\\0&1&0\\k&0&1\end{bmatrix}}_{\text{add }k\text{ row 1 to row 3}}.
\]
These are \textbf{elementary matrices}. They are old friends from Part 3: a swap
is a reflection (it exchanges two coordinates), a scaling is a non-uniform
scaling, and a replacement is a shear. Their determinants are $-1$, $c$ and $1$,
which are the laws D5, D4 and D8 of \S38. For example, with $k=-1$,
\[
  \begin{bmatrix}1&0&0\\0&1&0\\-1&0&1\end{bmatrix}
  \begin{bmatrix}1&2&0\\0&1&1\\1&0&1\end{bmatrix}
  =\begin{bmatrix}1&2&0\\0&1&1\\0&-2&1\end{bmatrix},
\]
which is ``subtract row 1 from row 3''.

%==============================================================================
\section{Triangular systems}
%==============================================================================

%------------------------------------------------------------------------------
\subsection{Back substitution}
%------------------------------------------------------------------------------
A triangular system is the easy case, and the goal of elimination is to reach
one. If $\mt{U}$ is upper triangular, the last equation of $\mt{U}\vc{x}=\vc{c}$
involves only $x_n$; the one above it, only $x_{n-1}$ and $x_n$; and so on. Solve
from the bottom up.

\begin{ex}
\label{ex:back}
\[
  \begin{aug}{ccc|c}2&1&-1&3\\0&3&2&7\\0&0&4&8\end{aug}
  \qquad
  \begin{aligned}
    x_3&=8/4=2,\\
    x_2&=(7-2x_3)/3=(7-4)/3=1,\\
    x_1&=(3-x_2+x_3)/2=(3-1+2)/2=2.
  \end{aligned}
\]
Check in the first equation: $2(2)+1-2=3$. \checkmark
\end{ex}

\noindent In general, once $x_{i+1},\dots,x_n$ are known, equation $i$ gives
\[
  x_i=\frac{1}{u_{ii}}\Bigl(c_i-\sum_{j=i+1}^{n}u_{ij}x_j\Bigr).
\]

\begin{algo}{Back substitution}{%
  \textbf{Input:} an upper triangular $\mt{U}\in\R^{n\times n}$ with every
  $u_{ii}\neq0$, and $\vc{c}\in\R^{n}$.\quad
  \textbf{Output:} the solution $\vc{x}$ of $\mt{U}\vc{x}=\vc{c}$.}
\item \kw{for} $i = n, n-1,\dots,1$:
\item \ind $s \;\gets\; c_i$
\item \ind \kw{for} $j = i+1,\dots,n$:
\item \indd $s \;\gets\; s-u_{ij}\,x_j$
        \hfill\textit{\small remove the unknowns already found}
\item \ind $x_i \;\gets\; s\,/\,u_{ii}$
\item \kw{return} $\vc{x}$
\end{algo}

%------------------------------------------------------------------------------
\subsection{Forward substitution}
%------------------------------------------------------------------------------
For a lower triangular $\mt{L}$, the first equation involves only $x_1$, so solve
from the top down:
$x_i=\frac{1}{l_{ii}}\bigl(c_i-\sum_{j<i}l_{ij}x_j\bigr)$ for $i=1,\dots,n$.

\begin{ex}
\[
  \begin{aug}{ccc|c}1&0&0&1\\2&1&0&4\\-1&3&2&7\end{aug}
  \qquad
  x_1=1,\quad x_2=4-2(1)=2,\quad x_3=\bigl(7+1-3(2)\bigr)/2=1 .
\]
\end{ex}

%------------------------------------------------------------------------------
\subsection{When it works, and what it costs}
%------------------------------------------------------------------------------
Substitution divides by every diagonal entry, so it needs them all nonzero. For a
triangular matrix that is exactly the condition $\det\mt{U}=u_{11}\cdots u_{nn}\neq0$
(\S39.1). If some $u_{ii}=0$, the system has either no solution or infinitely
many, and \S43 shows how to tell which.

Row $i$ costs about $n-i$ multiplications and as many subtractions, so the whole
substitution costs about $1+2+\cdots+n\approx\frac{n^{2}}{2}$ of each, roughly
$n^{2}$ operations. For $n=1000$ that is a million operations, a few
milliseconds of computing time. Triangular systems are cheap.

%==============================================================================
\section{Gaussian elimination and row-echelon form}
%==============================================================================

%------------------------------------------------------------------------------
\subsection{The idea}
%------------------------------------------------------------------------------
Use row operations, which do not change the solutions, to turn $[\,\mt{A}\mid\vc{b}\,]$
into $[\,\mt{U}\mid\vc{c}\,]$ with $\mt{U}$ upper triangular, then back
substitute. Part 1 \S4.4 hinted at this when it introduced triangular matrices.
The method works one column at a time: use the diagonal entry of the column, the
\textbf{pivot}, to clear every entry below it to zero.

%------------------------------------------------------------------------------
\subsection{Part 1's system, solved}
%------------------------------------------------------------------------------
\begin{ex}
\label{ex:part1}
Start from the augmented matrix of \S40.1.

\emph{Column 1.} The pivot is the $1$ in row 1. Row 2 already has a $0$ there.
Row 3 has a $1$, so subtract row 1 from row 3:
\[
  \begin{aug}{ccc|c}1&2&0&5\\0&1&1&5\\1&0&1&4\end{aug}
  \;\xrightarrow{R_3\gets R_3-R_1}\;
  \begin{aug}{ccc|c}1&2&0&5\\0&1&1&5\\0&-2&1&-1\end{aug}.
\]
\emph{Column 2.} The pivot is the $1$ in row 2. Clear the $-2$ below it by adding
twice row 2 to row 3:
\[
  \;\xrightarrow{R_3\gets R_3+2R_2}\;
  \begin{aug}{ccc|c}1&2&0&5\\0&1&1&5\\0&0&3&9\end{aug}.
\]
Upper triangular. Back substitute: $x_3=9/3=3$, $x_2=5-x_3=2$,
$x_1=5-2x_2=1$. So $\vc{x}=[\,1,2,3\,]$, the coordinates Part 1 \S1.6 built the
example from, now \emph{found} rather than known in advance.
\end{ex}

%------------------------------------------------------------------------------
\subsection{The algorithm}
%------------------------------------------------------------------------------
For an $n\times n$ system, the multiplier that clears entry $(i,k)$ is
$m_{ik}=a_{ik}/a_{kk}$, the entry divided by the pivot above it.

\begin{algo}{Gaussian elimination (forward phase)}{%
  \textbf{Input:} the augmented matrix $[\,\mt{A}\mid\vc{b}\,]$ with
  $\mt{A}\in\R^{n\times n}$.\quad
  \textbf{Output:} $[\,\mt{U}\mid\vc{c}\,]$ with $\mt{U}$ upper triangular and the
  same solutions.}
\item \kw{for} $k = 1,\dots,n-1$:
        \hfill\textit{\small work on column $k$}
\item \ind \kw{if} $a_{kk}=0$:
\item \indd \kw{if} a lower row has a nonzero entry in column $k$: swap it
        with row $k$
\item \indd \kw{else}: \kw{continue}
        \hfill\textit{\small no pivot in this column}
\item \ind \kw{for} $i = k+1,\dots,n$:
\item \indd $m \;\gets\; a_{ik}\,/\,a_{kk}$
\item \indd row $i$ $\;\gets\;$ row $i$ $-\;m\cdot$ row $k$
        \hfill\textit{\small now $a_{ik}=0$}
\item \kw{return} $[\,\mt{U}\mid\vc{c}\,]$
\end{algo}

\noindent Then back substitution finishes the job, if every diagonal entry of
$\mt{U}$ is nonzero. Counting operations: clearing column $k$ updates $n-k$ rows
of about $n-k$ entries each, about $2(n-k)^{2}$ operations, and summing over $k$
gives about $\frac23n^{3}$. Elimination is the expensive part. For $n=1000$ it is
under a billion operations, which a laptop does in about a second.

%------------------------------------------------------------------------------
\subsection{Zero pivots, and small ones}
%------------------------------------------------------------------------------
\begin{ex}
$\begin{aug}{cc|c}0&1&2\\1&1&3\end{aug}$: the first pivot position holds a $0$, and
nothing can be divided by it. Line 3 swaps the rows:
$\begin{aug}{cc|c}1&1&3\\0&1&2\end{aug}$, already triangular, so $x_2=2$ and
$x_1=1$.
\end{ex}

\begin{rmk}
On a computer, a pivot that is merely \emph{small} is nearly as bad as a zero
one: dividing by it magnifies every rounding error in its row. So practical codes
swap at every column, bringing up whichever entry at or below the diagonal is
largest in absolute value. This is \emph{partial pivoting}, and it is how the
algorithm is actually run.
\end{rmk}

%------------------------------------------------------------------------------
\subsection{Row-echelon form}
%------------------------------------------------------------------------------
Elimination works on any matrix, square or not, singular or not, if it is allowed
to move right whenever a column has no pivot. What it produces has a definite
shape.

\begin{defn}[Row-echelon form]
A matrix is in \textbf{row-echelon form} if
\begin{enumerate}[label=(\arabic*),leftmargin=2em,itemsep=2pt]
  \item any rows of all zeros are at the bottom, and
  \item the first nonzero entry of each nonzero row (its \textbf{pivot}) lies
        strictly to the right of the pivot of the row above.
\end{enumerate}
The columns that contain pivots are the \textbf{pivot columns}.
\end{defn}

\noindent The name is a picture: the pivots descend in steps, like an echelon of
troops. ($\ast$ marks any number; boxes mark pivots.)
\[
  \begin{bmatrix}\boxed{2}&\ast&\ast&\ast\\0&\boxed{1}&\ast&\ast\\0&0&0&\boxed{5}\\0&0&0&0\end{bmatrix}
  \text{ is in row-echelon form;}
  \qquad
  \begin{bmatrix}1&2\\3&4\end{bmatrix},\ \begin{bmatrix}0&1\\1&0\end{bmatrix}
  \text{ are not.}
\]
For a square matrix, row-echelon form is upper triangular.

%------------------------------------------------------------------------------
\subsection{The determinant, by elimination}
%------------------------------------------------------------------------------
For a square $\mt{A}$, elimination uses only swaps (each changes the sign of
the determinant, D5) and replacements (no change, D8), and ends at an upper
triangular $\mt{U}$ whose determinant is the product of its diagonal (\S39.1). So
\[
  \boxed{\;\det\mt{A}=(-1)^{s}\,u_{11}u_{22}\cdots u_{nn},\;}
\]
where $s$ is the number of swaps. In Example~\ref{ex:part1} there were no swaps
and the diagonal was $1,1,3$, so $\det\mt{A}=3$. A cofactor expansion agrees:
$1(1-0)-2(0-1)+0=3$. At about $\frac23n^{3}$ operations, this is how
determinants are computed in practice, not by the $n!$ of cofactors.

This also completes law D10 of \S38. Row operations do not change solutions, so
$\nullsp{\mt{A}}=\nullsp{\mt{U}}$. If every $u_{ii}\neq0$, back substitution on
$\mt{U}\vc{x}=\vc{0}$ gives only $\vc{x}=\vc{0}$. If some $u_{kk}=0$ (take the
first), the first $k$ columns of $\mt{U}$ have nonzero entries only in the first
$k-1$ rows, and $k$ vectors confined to $k-1$ coordinates are dependent. So
\[
  \nullsp{\mt{A}}=\{\vc{0}\}
  \iff \text{every }u_{ii}\neq0
  \iff \det\mt{A}\neq0,
\]
and with rank--nullity, $\rank\mt{A}=n\iff\det\mt{A}\neq0$.

%==============================================================================
\section{Gauss--Jordan elimination and reduced row-echelon form}
%==============================================================================

%------------------------------------------------------------------------------
\subsection{Going further}
%------------------------------------------------------------------------------
Gaussian elimination stops at triangular and hands over to back substitution.
\textbf{Gauss--Jordan elimination} keeps going with row operations instead: it
scales every pivot to $1$ and clears the entries \emph{above} each pivot as well
as below.

\begin{defn}[Reduced row-echelon form]
A matrix is in \textbf{reduced row-echelon form} (RREF) if it is in row-echelon
form, every pivot equals $1$, and every pivot is the only nonzero entry in its
column.
\end{defn}

\noindent Every matrix can be brought to reduced row-echelon form, and (we state
without proof) the result does not depend on which row operations were used: a
matrix has exactly one RREF.

%------------------------------------------------------------------------------
\subsection{A square system}
%------------------------------------------------------------------------------
\begin{ex}
Continue Example~\ref{ex:part1} from its triangular form:
\[
  \begin{aligned}
  &\begin{aug}{ccc|c}1&2&0&5\\0&1&1&5\\0&0&3&9\end{aug}
  \xrightarrow{R_3\gets\frac13R_3}
  \begin{aug}{ccc|c}1&2&0&5\\0&1&1&5\\0&0&1&3\end{aug}\\[4pt]
  &\qquad\xrightarrow{R_2\gets R_2-R_3}
  \begin{aug}{ccc|c}1&2&0&5\\0&1&0&2\\0&0&1&3\end{aug}
  \xrightarrow{R_1\gets R_1-2R_2}
  \begin{aug}{ccc|c}1&0&0&1\\0&1&0&2\\0&0&1&3\end{aug}.
  \end{aligned}
\]
The left block is now $\mt{I}$, and the system reads $x_1=1$, $x_2=2$, $x_3=3$
directly.
\end{ex}

\noindent For a square system with a unique solution, Gauss--Jordan ends at
$[\,\mt{I}\mid\vc{x}\,]$. Its real advantage shows when the solution is
\emph{not} unique.

%------------------------------------------------------------------------------
\subsection{Pivot variables and free variables}
%------------------------------------------------------------------------------
\begin{ex}
\label{ex:free}
Three equations, four unknowns:
\[
  \begin{aug}{cccc|c}1&2&0&1&1\\2&4&1&4&4\\3&6&1&5&5\end{aug}
  \xrightarrow[R_3\gets R_3-3R_1]{R_2\gets R_2-2R_1}
  \begin{aug}{cccc|c}1&2&0&1&1\\0&0&1&2&2\\0&0&1&2&2\end{aug}
  \xrightarrow{R_3\gets R_3-R_2}
  \begin{aug}{cccc|c}1&2&0&1&1\\0&0&1&2&2\\0&0&0&0&0\end{aug}.
\]
Column 2 had no pivot (after column 1 was cleared, its entries below row 1 were
all $0$), so the second pivot is in column 3. The result is already in reduced
row-echelon form, and it says
\[
  x_1+2x_2+x_4=1,\qquad x_3+2x_4=2,\qquad 0=0 .
\]
The pivot columns are $1$ and $3$; their unknowns $x_1,x_3$ are \textbf{pivot
variables}. The others, $x_2$ and $x_4$, are \textbf{free variables}: nothing
constrains them. Set $x_2=s$ and $x_4=t$, and solve for the pivot variables:
$x_1=1-2s-t$, $x_3=2-2t$. Every solution is
\[
  \vc{x}=
  \underbrace{\begin{bmatrix}1\\0\\2\\0\end{bmatrix}}_{\vc{x}_p}
  +s\begin{bmatrix}-2\\1\\0\\0\end{bmatrix}
  +t\begin{bmatrix}-1\\0\\-2\\1\end{bmatrix}.
\]
This is the shape promised in \S31.6: one solution $\vc{x}_p$ (take $s=t=0$) plus
the null space. Check: $\mt{A}\vc{x}_p=[\,1,\ 2+2,\ 3+2\,]=[\,1,4,5\,]$, and each
of the two other vectors is sent to $\vc{0}$.
\end{ex}

%------------------------------------------------------------------------------
\subsection{Reading off rank, nullity and bases}
%------------------------------------------------------------------------------
Row operations do not change which combinations of the columns are zero, because
$\mt{A}\vc{x}=\vc{0}$ and $\mt{R}\vc{x}=\vc{0}$ have the same solutions when $\mt{R}$ is
the RREF of $\mt{A}$. In $\mt{R}$, the pivot columns are distinct standard basis
vectors, so they are independent, and each non-pivot column is a combination of the
pivot columns to its left (its entries are the coefficients). The same relations
hold among the columns of $\mt{A}$. Hence:

\begin{thm}
\label{thm:pivots}
For any matrix $\mt{A}$, with $\mt{R}$ its reduced row-echelon form:
\begin{enumerate}[label=(\alph*),leftmargin=2em,itemsep=2pt]
  \item the columns of $\mt{A}$ in the pivot positions form a basis of
        $\colsp{\mt{A}}$, so $\rank\mt{A}$ is the number of pivots;
  \item setting one free variable to $1$ and the others to $0$, for each free
        variable in turn, gives a basis of $\nullsp{\mt{A}}$, so
        $\nullity\mt{A}$ is the number of free variables.
\end{enumerate}
\end{thm}

\noindent For (b), the vectors are independent because each has a $1$ in its own
free position, where all the others have $0$.

Every column is either a pivot column or a free column, so
\[
  \underbrace{\#\text{pivots}}_{\rank}+\underbrace{\#\text{free}}_{\nullity}=n .
\]
This is the rank--nullity theorem of \S32 again, now proved by counting columns.
In Example~\ref{ex:free}, $\rank=2$ (a basis of the column space is
$[\,1,2,3\,]$ and $[\,0,1,1\,]$, columns $1$ and $3$ of $\mt{A}$), $\nullity=2$,
and $n=4$.

%==============================================================================
\section{Homogeneous systems}
%==============================================================================

%------------------------------------------------------------------------------
\subsection{Always consistent}
%------------------------------------------------------------------------------
A homogeneous system $\mt{A}\vc{x}=\vc{0}$ always has the \textbf{trivial
solution} $\vc{x}=\vc{0}$. Its augmented column is zero and stays zero under every
row operation, so in practice one just reduces $\mt{A}$. The only question is
whether there are \textbf{nontrivial} solutions as well, and by
Theorem~\ref{thm:pivots} the answer is a count:
\[
  \text{nontrivial solutions exist}
  \iff \text{there is a free variable}
  \iff \nullity\mt{A}>0
  \iff \rank\mt{A}<n .
\]

\begin{thm}
A homogeneous system with more unknowns than equations ($m<n$) always has
nontrivial solutions.
\end{thm}

\begin{proof}
There is at most one pivot per row, so at most $m<n$ pivots, and at least one
column is free.
\end{proof}

%------------------------------------------------------------------------------
\subsection{An example}
%------------------------------------------------------------------------------
\begin{ex}
$x+y+z=0$ and $x-y=0$:
\[
  \begin{bmatrix}1&1&1\\1&-1&0\end{bmatrix}
  \xrightarrow{R_2\gets R_2-R_1}
  \begin{bmatrix}1&1&1\\0&-2&-1\end{bmatrix}
  \xrightarrow{R_2\gets-\frac12R_2}
  \begin{bmatrix}1&1&1\\0&1&\frac12\end{bmatrix}
  \xrightarrow{R_1\gets R_1-R_2}
  \begin{bmatrix}1&0&\frac12\\0&1&\frac12\end{bmatrix}.
\]
$z$ is free; $x=-\frac12z$ and $y=-\frac12z$. With $z=2t$, the solutions are
$t[\,-1,-1,2\,]$, a line in $\R^{3}$. Check: $-1-1+2=0$ and $-1+1=0$.
\end{ex}

%------------------------------------------------------------------------------
\subsection{When a system is consistent}
%------------------------------------------------------------------------------
For a general $\vc{b}$, elimination on $[\,\mt{A}\mid\vc{b}\,]$ settles everything.
A system is \textbf{consistent} if it has at least one solution.
\begin{itemize}[leftmargin=1.4em,itemsep=3pt]
  \item If the row-echelon form has a row $[\,0\ \cdots\ 0\mid c\,]$ with
        $c\neq0$, that row says $0=c$, and there is \textbf{no solution}.
        Otherwise the system is consistent.
  \item Consistent with no free variable: \textbf{exactly one} solution.
  \item Consistent with a free variable: \textbf{infinitely many}.
\end{itemize}
A row $[\,0\ \cdots\ 0\mid c\,]$ with $c\neq0$ means the last column of
$[\,\mt{A}\mid\vc{b}\,]$ is a pivot column. Equivalently, appending $\vc{b}$
raised the rank, or $\vc{b}\notin\colsp{\mt{A}}$:
\[
  \mt{A}\vc{x}=\vc{b}\text{ is consistent}
  \iff\rank[\,\mt{A}\mid\vc{b}\,]=\rank\mt{A}
  \iff\vc{b}\in\colsp{\mt{A}} .
\]

\begin{ex}
\S2.6 P5, one last time:
$\begin{aug}{cc|c}1&2&3\\2&4&1\end{aug}\xrightarrow{R_2\gets R_2-2R_1}\begin{aug}{cc|c}1&2&3\\0&0&-5\end{aug}$.
The second row says $0=-5$: no solution, found by the algorithm without any
picture.
\end{ex}

%==============================================================================
\section{Overdetermined and underdetermined systems}
%==============================================================================

%------------------------------------------------------------------------------
\subsection{More equations than unknowns}
%------------------------------------------------------------------------------
A system with $m>n$ is \textbf{overdetermined}. Its column space has dimension at
most $n<m$, a thin subspace of $\R^{m}$, and a randomly chosen $\vc{b}$ almost
never lies in it. Overdetermined systems are usually inconsistent, and they are
what measured data produce.

\begin{ex}
\label{ex:linefit}
Find a line $y=c_0+c_1t$ through the three data points $(1,1)$, $(2,2)$,
$(3,2)$. Each point gives one equation in the two unknowns $c_0,c_1$:
\[
  \begin{aug}{cc|c}1&1&1\\1&2&2\\1&3&2\end{aug}
  \xrightarrow[R_3\gets R_3-R_1]{R_2\gets R_2-R_1}
  \begin{aug}{cc|c}1&1&1\\0&1&1\\0&2&1\end{aug}
  \xrightarrow{R_3\gets R_3-2R_2}
  \begin{aug}{cc|c}1&1&1\\0&1&1\\0&0&-1\end{aug}.
\]
The last row says $0=-1$. No line passes through all three points, which is
obvious once they are plotted. Yet there is clearly a line that passes
\emph{close} to all three. Finding it is the subject of Part 9, which uses these
same three points.
\end{ex}

%------------------------------------------------------------------------------
\subsection{Fewer equations than unknowns}
%------------------------------------------------------------------------------
A system with $m<n$ is \textbf{underdetermined}. By \S44.1 the null space is
nonzero, with nullity at least $n-m$, so if the system is consistent it has
infinitely many solutions.

\begin{ex}
$x_1+x_2+x_3=3$: one equation, three unknowns, a whole plane of solutions,
$[\,3,0,0\,]+s[\,-1,1,0\,]+t[\,-1,0,1\,]$. Among them, $[\,3,0,0\,]$ has length
$3$ and $[\,1,1,1\,]$ has length $\sqrt3$. Is there a best one? A natural choice is
the \emph{shortest}, and \S54 shows it is $[\,1,1,1\,]$, the solution lying
along the row $[\,1,1,1\,]$ itself.
\end{ex}

%------------------------------------------------------------------------------
\subsection{Summary}
%------------------------------------------------------------------------------
\begin{center}
\renewcommand{\arraystretch}{1.35}
\begin{tabular}{@{}p{3.6cm} p{4.6cm} p{6.4cm}@{}}
\toprule
\textbf{Shape} & \textbf{Typical behaviour} & \textbf{What decides it} \\
\midrule
square, $\rank=n$ & exactly one solution, for every $\vc{b}$ & $\det\mt{A}\neq0$ \\
square, $\rank<n$ & none, or infinitely many & whether $\vc{b}\in\colsp{\mt{A}}$ \\
overdetermined, $m>n$ & usually none & whether $\vc{b}\in\colsp{\mt{A}}$; else least squares (Part 9) \\
underdetermined, $m<n$ & none, or infinitely many & if consistent, choose the shortest (\S54) \\
\bottomrule
\end{tabular}
\end{center}

%------------------------------------------------------------------------------
\subsection{Practice}
%------------------------------------------------------------------------------
\begin{enumerate}[label=\textbf{P\arabic*.},leftmargin=2.6em,itemsep=5pt]
  \item Solve by Gaussian elimination and back substitution:
        $x+y+z=6$, $2x+3y+z=11$, $x-y+2z=5$. Read off the determinant of the
        coefficient matrix from the pivots.
  \item Bring $\begin{aug}{ccc|c}1&2&1&3\\2&4&3&7\end{aug}$ to reduced row-echelon
        form and write every solution as $\vc{x}_p$ plus the null space. Give the
        rank and nullity.
  \item Is $x+y=2$, $x-y=0$, $2x+y=3$ consistent? If so, how many solutions?
  \item Write the elementary matrix that adds $5$ times row 1 to row 3 of a
        $3\times3$ matrix, and the elementary matrix that undoes it.
  \item Compute $\det\begin{bmatrix}0&2&1\\1&1&1\\2&0&3\end{bmatrix}$ by
        elimination, and check by cofactors.
  \item Without computing, how many free variables at least does a homogeneous
        system of $3$ equations in $5$ unknowns have?
\end{enumerate}

\subsubsection*{Answers}

\noindent\textbf{P1.} $R_2-2R_1$: $[\,0,1,-1\mid-1\,]$; $R_3-R_1$:
$[\,0,-2,1\mid-1\,]$; $R_3+2R_2$: $[\,0,0,-1\mid-3\,]$. So $z=3$, $y=-1+3=2$,
$x=6-2-3=1$. No swaps and pivots $1,1,-1$, so $\det=-1$.

\smallskip
\noindent\textbf{P2.} $R_2-2R_1$: $[\,0,0,1\mid1\,]$; $R_1-R_2$:
$[\,1,2,0\mid2\,]$. RREF
$\begin{aug}{ccc|c}1&2&0&2\\0&0&1&1\end{aug}$; $x_2=s$ free, $x_1=2-2s$, $x_3=1$.
$\vc{x}=[\,2,0,1\,]+s[\,-2,1,0\,]$. Rank $2$, nullity $1$.

\smallskip
\noindent\textbf{P3.} $R_2-R_1$: $[\,0,-2\mid-2\,]$; $R_3-2R_1$:
$[\,0,-1\mid-1\,]$; $R_3-\frac12R_2$: $[\,0,0\mid0\,]$. Consistent, no free
variable: exactly one solution, $y=1$, $x=1$. An overdetermined system
\emph{can} be consistent; this one happens to be.

\smallskip
\noindent\textbf{P4.} $\begin{bmatrix}1&0&0\\0&1&0\\5&0&1\end{bmatrix}$, undone by
$\begin{bmatrix}1&0&0\\0&1&0\\-5&0&1\end{bmatrix}$.

\smallskip
\noindent\textbf{P5.} Swap $R_1,R_2$: $\begin{bmatrix}1&1&1\\0&2&1\\2&0&3\end{bmatrix}$;
$R_3-2R_1$: $[\,0,-2,1\,]$; $R_3+R_2$: $[\,0,0,2\,]$. Pivots $1,2,2$, one swap:
$\det=-4$. Cofactors along row 1: $0-2(3-2)+1(0-2)=-4$. \checkmark

\smallskip
\noindent\textbf{P6.} At most $3$ pivots among $5$ columns, so at least $2$ free
variables.

\lookahead Example~\ref{ex:linefit} ended with no solution and an obvious best
approximate one. Part 9 makes ``best'' precise: the $\vc{x}$ that makes
$\mt{A}\vc{x}$ as close as possible to $\vc{b}$. It finds that $\vc{x}$ by solving a
different, always-consistent square system, using the elimination of this part.

\end{document}
